\section{Stochastic modelling}

Generally, at the outset of an infection, a relatively small number of colony
forming units --- bacteria or virions --- will enter the body. Be it during
inhalation into the respiratory system, through the skin with an insect's bite,
or via ingestion, those initial pathogenic particles will be critical in
determining the fate of the infection.

The founding population will either be depleted to extinction, or expand to a
capacity where its growth becomes established and requires a more targeted
effort of the immune system to deal with. Which of the two happens depends on
causes of a chance nature: encounters with cells and other immune factors that
are subject to the randomness of Brownian motion.

Stochastic effects play a more significant role in a population which is small.
As the count of individuals increases, the dynamics of overall growth and
decline become increasingly subject to the law of large numbers, and the
expected fate of a particular bacterium becomes telling of the group's
trajectory. Below that point a deterministic rate equation answers the wrong
question. It returns a single trajectory, growing or decaying, when what one
wants to know is whether the population establishes itself at all, and how
large it is if it does. Both are questions about a distribution, and neither is
visible in a mean.

For this reason, because we are interested in the early stages of a \yp{}
infection, it will be important to use stochastic models in our analysis.
